\documentclass[a4paper, landscape, 10pt]{article}
\usepackage[utf8]{inputenc}
\usepackage[T1]{fontenc}
\usepackage[ngerman]{babel}
\usepackage[left=1cm, right=1cm, top=1.5cm, bottom=1.5cm]{geometry}
\usepackage{helvet}
\usepackage{xcolor}
\usepackage{colortbl}
\usepackage{array}
\usepackage{ragged2e}

% Schriftart und Layout
\renewcommand{\familydefault}{\sfdefault}
\renewcommand{\arraystretch}{1.5}
\setlength{\parindent}{0pt}
\setlength{\parskip}{0pt}

% Farben
\definecolor{headergray}{RGB}{230, 230, 230}
\definecolor{textgray}{RGB}{50, 50, 50}

% Hilfsbefehl für Box-Titel (Robust für Tabellen)
\newcommand{\boxtitle}[1]{\textbf{\uppercase{#1}}}

% Spaltenbreiten vordefinieren
% Wir teilen die Seite in 5 logische Spalten
% Gesamtbreite ca 27.5cm. Eine Spalte ca 0.19\textwidth
\newlength{\colW}
\setlength{\colW}{0.192\textwidth}

\begin{document}

% --------------------------------------------------------------------------
% KOPFZEILE
% --------------------------------------------------------------------------
\begin{minipage}{0.7\textwidth}
    {\huge \textbf{THE PROJECT CANVAS}} \\
    % HIER WAR DER FEHLER: " durch \glqq und \grqq ersetzt
    {\large Masterarbeit: Generative \glqq Architecture is Code\grqq}
\end{minipage}
\hfill
\begin{minipage}{0.25\textwidth}
    \raggedleft
    \textbf{Erstellt von:} Moritz Diez \\
    \textbf{Datum:} \today
\end{minipage}

\vspace{0.5cm}

% --------------------------------------------------------------------------
% TABELLE
% --------------------------------------------------------------------------

% Struktur: 5 Spalten
\begin{tabular}{|p{\colW}|p{\colW}|p{\colW}|p{\colW}|p{\colW}|}
    \hline
    
    % ZEILE 1: ZWECK (2 Spalten) & KUNDE (3 Spalten)
    \multicolumn{2}{|l|}{\cellcolor{headergray}\boxtitle{ZWECK (Purpose)}} & 
    \multicolumn{3}{l|}{\cellcolor{headergray}\boxtitle{KUNDE (Customer)}} \\
    
    \multicolumn{2}{|p{0.384\textwidth}|}{
        \footnotesize \RaggedRight
        \textbf{Mehrwert \& Wissenszuwachs:} \par
        Das \glqq Schuhschachtel-Dilemma\grqq{} (unstrukturierte Altdaten) wird durch einen KI-Ansatz gelöst. Es wird aufgezeigt, wie SysML v2 als Enabler für generative Architektur-Synthese fungiert. \par
        \vspace{0.5em}
        \textbf{Fragen, die beantwortet werden:} \par
        Wie können valide Modelle aus Fragmenten generiert werden? Wie wird Modularität (\glqq Plug \& Play\grqq) ohne Refaktorierung sichergestellt?
    } & 
    \multicolumn{3}{p{0.576\textwidth}|}{
        \footnotesize \RaggedRight
        \textbf{Nutzer der Ergebnisse:} \par
        Unternehmen im Brownfield-Engineering; Systems Engineers, die vor der MBSE-Transformation stehen. \par
        \vspace{0.5em}
        \textbf{Anforderungssteller:} \par
        Hochschule München (Wissenschaftlicher Anspruch); Ggf. Industriepartner (Praxisrelevanz).
    } \\ \hline

    % ZEILE 2: ERGEBNIS (2 Spalten) & QUALITÄT (3 Spalten)
    \multicolumn{2}{|l|}{\cellcolor{headergray}\boxtitle{ERGEBNIS (Deliverables)}} & 
    \multicolumn{3}{l|}{\cellcolor{headergray}\boxtitle{QUALITÄT (Quality)}} \\
    
    \multicolumn{2}{|p{0.384\textwidth}|}{
        \footnotesize \RaggedRight
        \textbf{Objekte, die entstehen:} \par
        \textbullet\ \textbf{Masterarbeit (PDF):} Wissenschaftliche Dokumentation. \par
        \textbullet\ \textbf{AiC-Generator (Code):} Prototypische Implementierung der Pipeline (Ingestion, Matching, Synthesis). \par
        \textbullet\ \textbf{Modell-Beispiele:} Generierte SysML v2 Dateien (Teilmodelle \& Root-Modell).
    } & 
    \multicolumn{3}{p{0.576\textwidth}|}{
        \footnotesize \RaggedRight
        \textbf{Erfolgsmaßstäbe (Lösungsunabhängig):} \par
        \textbullet\ \textbf{Syntaktische Validität:} Der generierte Code muss vom SysML v2 Kernel fehlerfrei geparst werden. \par
        \textbullet\ \textbf{Modularität:} Teilmodelle müssen ohne manuelle Code-Änderung integrierbar sein. \par
        \textbullet\ \textbf{Reproduzierbarkeit:} Die Methodik muss auf unterschiedliche Datensätze anwendbar sein.
    } \\ \hline

    % ZEILE 3: MEILENSTEINE (3 Spalten) & RISIKEN (2 Spalten)
    \multicolumn{3}{|l|}{\cellcolor{headergray}\boxtitle{MEILENSTEINE}} & 
    \multicolumn{2}{l|}{\cellcolor{headergray}\boxtitle{RISIKEN + CHANCEN}} \\
    
    \multicolumn{3}{|p{0.576\textwidth}|}{
        \footnotesize \RaggedRight
        \textbf{Vorgehensmodell:} Inkrementelles Prototyping. \par
        \textbf{Wichtige Zeitpunkte:} \par
        \textbullet\ \textbf{15.03.26:} Projektstart / Kick-Off. \par
        \textbullet\ \textbf{15.04.26:} Konzept steht (Anforderungs-Review). \par
        \textbullet\ \textbf{31.05.26:} Generator-Prototyp läuft (\glqq Hello World\grqq{} der Generierung). \par
        \textbullet\ \textbf{15.07.26:} Validierung abgeschlossen (Integrationstest erfolgreich). \par
        \textbullet\ \textbf{15.09.26:} Abgabe der Arbeit.
    } & 
    \multicolumn{2}{p{0.384\textwidth}|}{
        \footnotesize \RaggedRight
        \textbf{Risiken (Scheitern):} \par
        \textbullet\ KI-Halluzinationen führen zu falschen Interfaces. \par
        \textbullet\ Komplexität der SysML v2 Grammatik wird unterschätzt. \par
        \textbf{Chancen:} \par
        \textbullet\ Blaupause für automatisierte Digitalisierung. \par
        \textbullet\ Beitrag zum Stand der Technik (\glqq AI for SE\grqq).
    } \\ \hline
    
    % ZEILE 4: RESSOURCEN (1), TEAM (2), BUDGET (2)
    \multicolumn{1}{|l|}{\cellcolor{headergray}\boxtitle{RESSOURCEN}} & 
    \multicolumn{2}{l|}{\cellcolor{headergray}\boxtitle{TEAM}} & 
    \multicolumn{2}{l|}{\cellcolor{headergray}\boxtitle{BUDGET}} \\
    
    \multicolumn{1}{|p{\colW}|}{
        \footnotesize \RaggedRight
        \textbf{Sachmittel:} \par
        \textbullet\ LLM-API Zugang. \par
        \textbullet\ SysML v2 Kernel. \par
        \textbullet\ Testdaten.
    } & 
    \multicolumn{2}{p{0.384\textwidth}|}{
        \footnotesize \RaggedRight
        \textbf{Beteiligte:} \par
        \textbullet\ Moritz Diez (Umsetzung). \par
        \textbullet\  Prof. XXX (Betreuer HM). \par
        \textbullet\ Kai Heinrich (Firmenbetreuer Rücker + Schindele).
    } & 
    \multicolumn{2}{p{0.384\textwidth}|}{
        \footnotesize \RaggedRight
        \textbf{Finanzmittel:} Keine. \par
        \textbf{Zeitbudget:} \par
        30 ECTS $\approx$ 900h Aufwand.
    } \\ \hline

    % ZEILE 5: UMFELD (2) & ZEIT (3)
    \multicolumn{2}{|l|}{\cellcolor{headergray}\boxtitle{UMFELD}} & 
    \multicolumn{3}{l|}{\cellcolor{headergray}\boxtitle{ZEIT (Time)}} \\
    
    \multicolumn{2}{|p{0.384\textwidth}|}{
        \footnotesize \RaggedRight
        \textbf{Rahmenbedingungen:} \par
        Masterstudiengang Systems Engineering (HM). \par
        Technologischer Wandel zu \glqq Everything as Code\grqq. \par
        Verfügbarkeit von LLMs.
    } & 
    \multicolumn{3}{p{0.576\textwidth}|}{
        \footnotesize \RaggedRight
        \textbf{Start Vorbereitung:} Februar 2026 \par
        \textbf{Start Bearbeitung (Offiziell):} 15.03.2026 \par
        \textbf{Abgabe:} 15.09.2026 (Fixer Termin)
    } \\ \hline

\end{tabular}

\vspace{0.5cm}
\tiny
\textit{Hinweis: PDF-Formular ist zwingend elektronisch auszufüllen (vgl. Infoblatt S. 7).}

\end{document}